\documentclass[10pt]{article}

\usepackage[utf8]{inputenc}

\usepackage[T1]{fontenc}

\usepackage{amsmath}

\usepackage{amsfonts}

\usepackage{amssymb}

\usepackage[version=4]{mhchem}

\usepackage{stmaryrd}



\begin{document}

мый термом $A$. Выражение (*) и наз. Ч. $\lambda$-а. Эта Ч. $\lambda$-а., наз. также явным определением функций, употребляется чаще всего в случае, когда в языке теории возникает опасность смешения функции как объекта исследования со значениями функции для нек-рых значений аргумента. Введена А. Чёрчем [1].



Jum.: [1] C h u r c h A., The calculi of lambda-conversion, Princeton, 1941; [2] К а p p и X. Б., Основания математической логики, пер. с англ., М., 1969. А. Г. Драгалин.



ЧЁРЧА ТЕЗИС - принцип, согласно к-рому класс функций, вычислимых с помощью алгоритмов в пироком интуитивном смысле, совпадает с классом частично рекурсивных функций. Ч. т. - это естественнонаучный факт, подтверждаемый опытом, накопленным в математике за всю ее историю. Все известные в математике примеры алгоритмов удовлетворяют ему. Ч. т. впервые был высказан А. Чёрчем (A. Church, 1936). Различным уточнениям интуитивного понятия алгоритма соответствуют свои формулировки Ч. т. Т е з и с $\mathbf{T}$ ь юр и н г а заключается в том, что всякая вычислимая в интуитивном смысле функция вычислима с помощью нек-рой Тьюринга машины, а принцип нормализации Маркова - в том, что всякая вычислимая в интуитивном смысле функция вычислима с помощью нек-рого нормального алгорифма. Из эквивалентности известных уточнений понятия алгоритма следует әквивалентность соответствующих вариантов Ч. т. Этот факт является еще одним подтверждением Ч. т. Тезис Чёрча не может быть строго доказан, так как в его формулировке участвует неточное понятие «алгоритм в интуитивном смысле». Были попытки опровергнуть Ч. т., однако они к успеху не привели (1984). Принятие Ч. т. полезно в теории алгоритмов и ее приложениях. Во-первых, при доказательстве существования тех или иных конкретных алгоритмов - машин Тьюринга, рекурсивных функций, нормальных алгорифмов и др. - можно, опираясь на Ч. т., ограничиваться интуитивно ясными построениями и не выписывать соответствующие формальные схемы.



Кроме того, Ч. т. является основанием для вывода о неразрешимости данной алгоритмической проблемы после того, как строго доказано, что эта проблема не может быть решена в рамках того или иного уточнения понятия алгоритма.



Лum.: [1] К ли н и С. К., Введение в метаматематику, пер. с англ., М., 1957; [2] Р о'д ж е р с X., Теория рекурсивных функций и әффективная вычислимость, пер. с англ., M., ных функций и әффективная вычислимость, пер. с англ., М.,

C. И. Аәян.



ЧЕТАЕВА ПРИНЩИП - вариационный диффферегциальный принцип механики, представляющий собой видоизменение Гаусса принципа, установленное Н. Г. Четаевым [1].



Согласно Ч. п., работа на элементарном ццкле, состоящем из прямого движения в поле заданных сил и движения обратного (попятного) в поле сил, к-рых было бы достаточно для создания действительного движения, если бы механич. система была совершенно свободной, для действительного движения имеет относительный (по меньшей мере) максимум в классе мыслимых по Гауссу движениӥ. Ч. п. распространен на физич. системы, а также на сплопшные среды (см. [2]).



Подробнее см. Вариационные принципы классической механики.



Лит.: [1] Ч е т а е в Н. Г., «Прикл. матем. и механ.», 1941, т. 5, в. 1, с. 11-12; [2] P у м я н ц е в B. В., "Докл. AH CCCP», 1973, т. $210, \mathfrak{N}_{2} 4$, с. 787-90. В. В. Румянчев.



ЧETAEBA ТЕOРЕМЫ-1) Ч. т. о н е у с т о й ч Ив о с т и - общие теоремы о неустойчивости движения, установленшые Н. Г. Четаевым для уравнений возмуценного движения вида



\[

\frac{d x_{s}}{d t}=X_{s}\left(t, x_{1}, \ldots, x_{n}\right), s=1, \ldots, n,

\]



правые части к-рых $X_{s}$ - голоморфные функции относительно действительных переменных $x_{s}$ с коэфффициентами, являющимися непрерывными и ограниченными функциями действительной переменной - времени $t$, определенные в нек-рой области



\[

t \geqslant t_{0}, \quad \sum_{s=1}^{n} x_{s}^{2}<A

\]



причем $X_{s}(t, 0, \ldots, 0)=0$.



Ч. т. даны в двух формулировках: с двумя функциями $V, W$ и с одной функцией $V$. Под функциями $V, W$ понимаются действительные функции действительных переменных $x_{s}$ и $t$, однозначные и непрерывные в области (2) и обращающиеся в нуль при $x_{s}=0$, так же как их полные производные по времени 1-го порядка $\dot{V}, \dot{W}$, причем, напр.,



\[

\dot{V}=\frac{\partial V}{\partial t}+\sum_{s=1}^{n} \frac{\partial V}{\partial x_{s}} X_{s}

\]



Те орема онеусто й ч ивостисдд у мя ф у н к ц и я м и (1934, см. [1], с. 222-24): если дифференциальные уравнения таковы, что: 1) для нек-рой допускающей бесконечно малый высший предел функции $V$ существует область, где $V \dot{V}>0$, и 2) для нек-рых значений величин $x_{s}$, численно сколь угодно малых, в области $V \dot{V}>0$ возможно выделить область, в к-рой нек-рая функция $W>0$, на границе области $W=0$ значения $\dot{W}$ суть одного какого-либо оюределенного знака, то невозмущенное движение $x=0$ неустойтиво.



При решении вопроса о неустойчивости целесообразно рассматривать ивтервал измешения времени $\left[t_{0}, \infty\right]$ закрытым и существование области $V>0$ понимать как eе непустоту для любого $t$ на этом интервале (см. [1], с 225-238). Если рассматриваемая область $V \dot{V}>0$ ограничена $V=0$ и при этом $\dot{V} \geqslant 0$, то за функцию $W$ возможно взять функцию $V$.



Условие 2) можно сформулировать многими иными способами, в частности, пусть уравнения (1) имеют первый интеграл $F(t, x)=$ const, обладающий свойствами функций $V$, и пусть область $F>0$ не является пустой при численно сколь угодно малых значениях $x_{s}$ (см. [1], с. 232). Если для возмущенных движений при $F(t, x)=\varepsilon$, сколь бы мало $\varepsilon$ не было, тем самым будут выполняться неравенства $\dot{V}>l^{\prime}>0$ в области (2) для допускающей малыӥ высший предел функции $V$, a начальные значения $x_{s_{0}}$ при удовлетворении интегралу $F(t, x)=\varepsilon=F_{0}$ возможно выбрать так, чтобы они удовлетворяли также и неравенству $V_{0}>0$, то невозмущенное движение неустойчиво.



Если интервалі времени считать открытым $\left[t_{0}, \infty\right)$ и не вводить условного смысла суцествования области $V>0$, то справедлива т е о р е м а о н е у с т о й ч ив о с т и с о д н о ї ф у н к и е ї $(1946, \mathrm{~cm}$. [1], с. 5-152): если диффференциальные уравнения (1) таковы, что возмоэкно найти функцию $V$, ограниченную в области $V>0$, существующей при всяком $t \geqslant t_{0}$ и для сколь угодно малых по абсолютной величине значений переменных $x_{s}$, ироизводная к-рой $\dot{V}$ была бы определенноположительной в области $V>0$, то невозмущенное движение неустойчиво. Год определенно положительной в области $V>0$ функцией понимается функция $W(t, x)$, к-рая может обращаться в нуль в әтой области лишь на границе $V=0$, причем для произвольного $\varepsilon>0$, как бы мало оно не было выбрано, найдется такое число $l>0$, что при $x_{s}$, удовлетворяющих условию $V \geqslant \varepsilon$, и для всякого $t \geqslant t_{0}$ имеет место неравенство $W \geqslant l$.



Было дано обращение этой теоремы, чем была обоснована ее универсальность (см. [2]).





\end{document}